\section{The solver module map}
\label{architecture-solver-map}

The execution machinery described later serves a set of physical and numerical
modules. Their boundaries are easiest to understand by following what each
consumes, what it produces, and which state must survive to the next call.
The following map covers the top-level solver directories under
\texttt{include/} and \texttt{src/}. It is an ownership index, not an assertion
that every factory entry or combination is Ready. The current catalog and
readiness registry remain authoritative for admission.

\begin{longtable}{@{}p{0.25\linewidth}p{0.70\linewidth}@{}}
\toprule
Module and interfaces & Implementation responsibility \\
\midrule
\endhead
\bottomrule
\endfoot
\href{../include/core}{Core interfaces} &
\href{../src/core}{Core implementation}: field storage, physical geometry,
configuration/runtime state, headless lifecycle, nested execution, checkpoint
and output owners; orchestration composes accepted operator changes. \\
\href{../include/init}{Initialization} &
\href{../src/init}{Typed configuration and initial sources}: sounding-column
construction, hodograph selection, trigger selection and requested/realized
environment diagnostics. \\
\href{../include/soundings}{Sounding readers} &
\href{../src/soundings}{File ingestion}: auto routing, CM1 and SHARPY-related
readers, unit conversion, profile quality control and interpolation. \\
\href{../include/dynamics}{Dynamics} &
\href{../src/dynamics}{Cartesian A-grid/C-grid tendencies}, reference-state
forces, native boundaries, conservative momentum and moving-frame owners. \\
\href{../include/numerics}{Numerics} &
\href{../src/numerics}{Advection, diffusion and integration}: TVD/WENO,
retained carrier fluxes, explicit/implicit diffusion, split-explicit/RK methods
and their stability contracts. \\
\href{../include/boundary_conditions}{Boundary conditions} &
\href{../src/boundary_conditions}{Geometric closure}: staggering-specific
boundary execution, closure storage and private transported-field boundaries. \\
\href{../include/microphysics}{Microphysics} &
\href{../src/microphysics}{Kessler, Lin, Thompson and Milbrandt}: water sources,
private particle moments, sedimentation, thermal attribution and state commit. \\
\href{../include/radiation}{Radiation} &
\href{../src/radiation}{Simple-grey and RRTMG}: physical columns, optics,
fluxes, heating, mapped shortwave transport and held-source energy accounting. \\
\href{../include/boundary_layer}{Boundary layer} &
\href{../src/boundary_layer}{Slab, YSU and MYNN}: column mixing, surface
exchange, persistent column state and conserved interval-mean forcing. \\
\href{../include/turbulence}{Turbulence} &
\href{../src/turbulence}{Resolved-grid SGS}: Smagorinsky variants and
prognostic TKE, coefficient fields, stress/scalar fluxes and production budgets. \\
\href{../include/terrain}{Terrain} &
\href{../src/terrain}{None/Bell/Sch\"ar topography}, coordinate metrics,
physical flux transforms and mapped transport geometry. \\
\href{../include/chaos}{Chaos} &
\href{../src/chaos}{Controlled stochastic perturbations}: keyed random
streams, spatial/temporal correlation, source coupling and retained epochs. \\
\href{../include/radar}{Radar} &
\href{../src/radar}{Synthetic observation operators}: reflectivity, radial
velocity and polarimetric products with declared geometry and particle inputs. \\
\href{../include/diagnostics}{Diagnostics} &
\href{../src/diagnostics}{Field contracts, budgets and derived products},
storm tracking, rotation diagnostics and parcel trajectories. \\
\href{../include/compute}{Compute} &
\href{../src/compute}{Backend-neutral runtime, CPU owners, typed dispatch
and residency}; native GPU implementations live under Vulkan and Metal. \\
\href{../include/util}{Utilities} &
Shared scheme registration, numeric text, logging, cadence, transient
workspace and copy helpers; relevant compiled helpers also live in
\href{../src/core/infra}{core infrastructure}. \\
\href{../src/tools}{Solver tools} &
Field-contract validation and Thompson activation/collection/cold-table
generation. These executables support provisioning and tests; they are not
additional atmospheric schemes. \\
\end{longtable}

There is no separate implemented soil or prognostic land-surface directory.
The similarly named \texttt{surface.scheme} and
\texttt{boundary\_layer.scheme} controls have different responsibilities,
explained below. A directory inventory must not invent a model because a
related name appears in configuration or in a paper.

Most physical families use the descriptor-backed
\href{../include/util/scheme_factory.hpp}{SchemeRegistry}. It validates
canonical IDs, aliases and constructor agreement. A constructor establishes
that an object can be built; it does not establish valid input geometry,
complete private-state continuation, native dispatch or a supported coupling.
Initialization sources and geometric boundaries have their own typed factory
dispatch. These different factories should not be forced into one lifecycle
merely because they all select an implementation.

\section{Initialization is a composition of three sources}
\label{architecture-initialization}

\subsection{Thermodynamic column construction and file reading are separate}

\texttt{src/init/sounding} owns the model-column source. Its choices are
\texttt{parametric\_cape}, \texttt{parametric\_targets} and a file-backed
source. The source returns a \texttt{Sounding} containing physical heights,
temperature, theta, dry-air vapor mixing ratio, pressure, dry density and
optional Cartesian winds. The caller broadcasts or physically samples that
column onto the selected grid. The source is not simply a filename reader.
See the \href{../include/init/sounding/sounding_source.hpp}{column interface}
and \href{../src/init/sounding/factory.cpp}{source factory}.

The separate \texttt{soundings} module supplies the auto, CM1 and SHARPY-related
file readers. It interprets text and supported netCDF/HDF5 representations,
checks required variables and units, and interpolates to requested heights.
Strict file initialization requires adequate vertical coverage and rejects
missing or invalid inputs rather than silently creating an unrelated profile.
Lower-level compatibility code contains explicit fallback options; their
existence does not change the ordinary strict-v1 no-fallback contract.

The \href{../src/init/sounding/sources/file_sounding.cpp}{file-column adapter}
reintegrates pressure on the model column. A file pressure at the first
interpolated scalar center is not automatically a surface-face pressure;
without that supplied center pressure, the surface boundary must be integrated
to the first center. Density and theta then follow the retained dry
thermodynamic convention. This separates observed input representation from the
reference state the discrete force operator actually consumes. It does not
prove exact hydrostatic balance for every possible later stencil; the
initialization and resting-column tests check the implemented pairing.

Parametric targets similarly produce a realized column and diagnostics.
Requesting CAPE, humidity or a characteristic level does not make every
environmental diagnostic independently controllable. Requested and realized
values must remain distinct. A parameter solver's failed or out-of-domain
target should not be reported as an exact physical initial condition just
because it returned a vector of the requested length.

\subsection{Wind selection and perturbations have explicit precedence}

The \href{../src/init/hodograph/factory.cpp}{hodograph factory} handles
\texttt{auto}, \texttt{wk\_param} and \texttt{zero}. Auto uses complete winds
from the sounding when available; otherwise it builds the retained parametric
hodograph. Explicit parametric or zero selection has its own meaning rather
than being silently replaced by file winds. The resulting Cartesian wind
column is placed on the chosen scalar/face locations by initialization
orchestration. Wind direction interpolation should operate through components
where required, avoiding a spurious angular jump across north.

The \href{../src/init/trigger/factory.cpp}{trigger factory} selects the warm
bubble, no perturbation or a vortex seed after the base column and wind state
exist. The current
\href{../src/core/orchestration/dynamics/initial_conditions_cartesian.cpp}{Cartesian
warm-bubble owner} is not theta-only: for positive heating with available
pressure/vapor fields, it rescales vapor using the ratio of its retained
saturation estimates at the old and new temperatures. Optional cloud-liquid
seeding is another declared increment. This is an imposed initial-state
perturbation, not a closed microphysical condensation transfer. Its moisture
and thermal assumptions must be included in reproducing an experiment.

The current vortex seed adds a horizontal Rankine profile about its x/y center
up to the hard \texttt{z\_top\_m} cutoff. Its public
\texttt{center\_z\_m} parameter is parsed and serialized but currently has no
effect in the seed application. The owner does not author a balancing pressure
field, and its discrete staggered field must not be described as a proven
balanced or divergence-free vortex. These are implementation limits to retain
when interpreting the \href{../src/init/trigger/sources/vortex_seed.cpp}{seed source}.

Terrain initialization uses physical sample heights, and subsequent reference
and EOS reconciliation belongs to orchestration. Trigger code should not be
treated as an independent complete thermodynamic integrator. The permanent
\href{../tests/init/test_file_sounding.cpp}{file-source},
\href{../tests/init/test_parametric_targets_sounding.cpp}{target-column},
\href{../tests/init/test_hodograph_sources.cpp}{hodograph} and
\href{../tests/init/test_vortex_seed_trigger.cpp}{trigger tests} define bounded
checks at these interfaces.

\section{Resolved dynamics, numerical operators and boundaries}

The \href{../src/dynamics/factory.cpp}{dynamics factory} constructs distinct
Cartesian A-grid and C-grid schemes. Their tendencies consume the state view,
base profiles and physical geometry; the orchestration layer owns applying
them in a valid sequence. Pressure gradients, gravity, vapor buoyancy,
condensate loading, rotation and advection are not interchangeable source
plugins with arbitrary ordering. A native operator must implement the selected
grid's actual field locations and force convention.

\texttt{numerics} owns the reusable spatial and temporal algorithms. Its
advection factory provides TVD, WENO5 and disabled transport; explicit and
implicit diffusion have separate operators and configured coefficient sources.
Time integration selects split-explicit, SSPRK3 or classical RK4. A disabled
configuration may bypass an operator rather than construct a meaningful
diffusion solve with zeroed input. Interface-flux registers retain the mass
carrier and interval needed by scalar transport and nested feedback.

The distinction between a tendency and a committed state matters here. An RK
stage can require primary fields, private microphysics state and reports to be
staged together; a split acoustic callback has different pressure/boundary
dependencies. The dry rho--theta EOS projection diagnoses pressure from stored
density and theta with its declared arithmetic. It does not become a moist
gas-law solve because a neighboring physics source includes water energy.
Numerical diffusion, SGS stress and PBL mixing also remain separate owners;
selecting one does not implicitly remove the others' coupling constraints.

\href{../src/boundary_conditions/factory.cpp}{Boundary construction} dispatches
on coordinate system and staggering, while its runtime chooses the level's
execution object. Physical walls, periodic connections and open/radiative
conditions define closure values and histories needed by later stencils.
They do not turn halo storage into physical atmospheric volume. Private
transported moments require corresponding boundary closure and cannot inherit
the primary water field's state by accident.

Surface exchange is not geometric boundary filling. A rigid lower normal-flow
condition can coexist with a surface stress or heat flux owned by the PBL
coupling. Conversely, copying a temperature into a ghost layer does not
automatically supply a physically accounted surface energy flux. Native
boundary, source and diagnostic owners must agree on that separation.
See \href{../src/boundary_conditions/runtime.cpp}{boundary runtime},
\href{../src/numerics/runtime.cpp}{numerics runtime} and
\href{../src/core/orchestration/dynamics/dynamics.cpp}{step orchestration}.

\section{Microphysics owns more than the visible water arrays}

The \href{../src/microphysics/factory.cpp}{microphysics registry} supplies
Kessler, Lin, Thompson, Milbrandt and disabled physics. The math guide describes
their different process and spectrum assumptions. Architecturally, the shared
interface accepts pressure, theta, live dry density, water species, physical
layer information and the interval, and produces primary changes plus any
scheme-owned moment changes and receipts.

Thompson and Milbrandt retain private fields that must follow transport,
boundary closure, cloning, relocation and checkpoint continuation. The
\href{../include/microphysics/microphysics_base.hpp}{base contract} distinguishes
parameter cloning from evolved state, exposes the storage units of transported
private scalars, and separates preparation/validation from publication.
An accepted q update with old N or Z would describe a different spectrum even
if total water still summed correctly.

\href{../src/core/orchestration/physics/microphysics_step.cpp}{The microphysics
step owner} integrates those results with primary-state acceptance and budgets.
Sedimentation is a physical lower-boundary exchange with a precipitation
receipt, not merely another local category conversion. Execution must reach
the process and commit owners declared by admission: an explicitly host-owned
commit is different from silently replacing a required device operation.
A device capability bit alone proves neither. Ordinary double-moment Milbrandt and the guarded
triple-moment work are separate scopes; the latter's bounded wet-surface/EOS
helpers do not complete public field-level admission.

\section{Radiation: columns, spectral transfer and held power}
\label{architecture-radiation}

\subsection{The column boundary makes physical inputs explicit}

The \href{../include/radiation/radiation_base.hpp}{radiation interface} defines
temperature, pressure, dry density, physical layer thicknesses and interface
heights, optional interface pressure/temperature, water profiles, gas
composition, solar geometry and surface optical/thermal state. Center pressure
is not an interchangeable substitute for pressure at the surface or top
interface. Likewise, nominal grid spacing cannot replace mapped column
thickness. \href{../src/radiation/column_evaluation.cpp}{Column evaluation}
checks dimensions, finite input, flux/tendency shape and the declared
column-power closure before publication.

The factory contains \texttt{none}, \texttt{simple\_grey} and \texttt{rrtmg}.
Omitted radiation means no radiative forcing, not an automatically inserted
clear-sky model. Simple-grey offers a pressure-based clear-sky optical profile
and a distinct absorbing vapor/liquid/ice-path model. The latter's pure
absorption assumption does not represent cloud scattering. The collimated
incident solar beam and Lambertian surface-reflected flux have different
angular meanings; sharing a broadband coefficient does not make them the
same ray distribution.

RRTMG is a C++ spectral implementation whose AER coefficient files are explicit
scientific inputs. Coefficient loading, atmosphere conversion, reference-grid
interpolation, gas/cloud optics and longwave/shortwave transfer have separate
source owners under \href{../src/radiation/schemes/rrtmg}{the RRTMG directory}.
It does not call a reference Fortran program at runtime. Its resolved-condensate
cloud optics represent a fully cloudy cell where applicable rather than
inventing a subgrid cloud-fraction/overlap model. Explicit gas mixing ratios
and particle effective radii are part of the closure, not hidden renderer
settings. Dry-gas profiles use dry-air mole/volume fractions, while water inputs
retain their declared dry-air mass ratios; the atmosphere adapter owns the
conversion to molecular columns. Missing coefficients or unsupported inputs
must not turn RRTMG into simple-grey.

\subsection{Heating and refresh cadence have separate ownership}

Radiative fluxes produce a column temperature tendency through the retained
physical mass and heat-capacity convention. That tendency is converted to
theta for coupling, while conserved power and flux diagnostics are published
with it. Radiation does not directly replace the dry EOS. Agreement of its
column power with integrated heating checks one source boundary; the later
thermal and pressure commits still need their own energy attribution.

\href{../src/radiation/runtime.cpp}{Radiation runtime} retains last/next sample
times and holds the accepted source between refreshes. It constructs candidate
column fields and reports, rejects failed columns before publishing the new
generation, and retains the previous accepted source when a refresh fails.
The cadence is part of continuation state. Recomputing an otherwise held source
immediately after restart can change subsequent evolution even if every visible
field initially matches.

The \href{../src/radiation/energy_ledger.cpp}{application ledger} compares
held atmospheric power times the applied interval against booked radiation
heat. Mapped shortwave also distinguishes incoming power, atmospheric and
surface absorption, top and lateral escape, and unresolved reflected power.
Reflection convergence is an acceptance condition, not a number to hide inside
the atmospheric tendency. These receipts describe actual source publication
and application; they do not by themselves prove the complete model conserves
a physical total energy.

\subsection{Terrain and native execution are additional contracts}

Terrain shortwave uses slope/visibility, physical interfaces, lateral boundary
semantics and angular/reflection work beyond independent vertical columns.
The current CPU mapped RRTMG path has a spectral terrain operator. Its radiation
runtime explicitly rejects mapped RRTMG shortwave on an active accelerator
because the required native spectral-terrain owner is missing. This limitation
must not be inferred away from the existence of RRTMG native column batches.

For supported column work, the runtime asks the scheme whether the selected
backend supplies the required native batch operations. Rejected or unavailable
dispatch fails a no-fallback profile; CPU fallback is used only when explicitly
permitted. Geometry and output-diagnostic combinations can impose further
restrictions. The term ``C++ implementation'' therefore does not mean every
spectral branch is device-native. See
\href{../src/radiation/native_batch_runtime.cpp}{batch publication},
\href{../tests/physics/test_radiation.cpp}{radiation references} and
\href{../tests/physics/test_radiation_runtime_state.cpp}{runtime-state checks}.

\section{SGS, boundary-layer and surface exchange are different models}
\label{architecture-mixing-surface}

\subsection{Resolved-grid SGS}

\texttt{turbulence} constructs Smagorinsky, prognostic TKE or disabled SGS.
Smagorinsky variants own coefficient/filter/stability choices; the TKE family
adds an evolved energy field, mixing length, shear and buoyancy production,
dissipation and transport. Their spatial stress and scalar fluxes must respect
the chosen staggering and physical metrics. Returning an eddy viscosity field
is not the complete operator: application to momentum/scalars, realized energy
change and boundary work remain separate steps.

\href{../src/turbulence/runtime.cpp}{SGS runtime} and
\href{../src/core/orchestration/dynamics/sgs_apply.cpp}{SGS application}
connect closure preparation to accepted changes. C-grid tensor placement,
collocated derivatives and mapped stresses have different contracts. Dynamic
Smagorinsky's bounded coefficient and averaging policy do not automatically
authorize a vertical-only variant; the retained static vertical-only scope is
still tracked separately as Engineering incomplete. A filter width is a
physical closure parameter even when its default is derived from the grid.

\subsection{Column PBL and interval-mean forcing}

\texttt{boundary\_layer} constructs slab, YSU, MYNN or disabled mixing. Slab
retains per-column mixed-layer depth and thermodynamic state; YSU and MYNN have
their own nonlocal/stability/TKE assumptions. Their column workspace receives
physical centers/interfaces and scalar-location wind reconstructed from the
actual grid. A PBL column is not a view of all allocated halo entries.

\href{../src/boundary_layer/column_evaluation.cpp}{Column evaluation}
checks completed results before
\href{../src/boundary_layer/runtime.cpp}{runtime publication}. The held
representation includes forces and scalar inventory rates, not just copied
accelerations whose denominator may later change. Applying a retained force
to the current face mass requires the appropriate finite-step momentum/energy
bookkeeping. Implicit column exchange publishes its realized interval-mean
fluxes; requested diffusivities alone do not describe that accepted transfer.
Persistent column state, cadence and held forcing belong in rollback and
restart.

Native batch submission is another owner, with bounded packing and complete
publication. Its presence does not license all cloud-top, terrain or SGS/PBL
compositions. Admission must prevent two closures from unintentionally modeling
the same unresolved transport twice. See
\href{../tests/physics/test_boundary_layer.cpp}{PBL reference checks} and
\href{../tests/physics/test_boundary_layer_runtime_state.cpp}{PBL state checks}.

\subsection{Surface state does not imply a soil model}

The public surface choices are \texttt{none} and \texttt{slab}. Here
surface selection controls lower-boundary exchange with an active PBL closure;
it is distinct from selecting the slab PBL family. Monin--Obukhov similarity
and related surface exchange live in the PBL implementation, including
\href{../src/boundary_layer/base/surface_fluxes.cpp}{surface-flux helpers}.
\texttt{SurfaceConfig} carries roughness, albedo, emissivity, skin temperature
and surface vapor mixing ratio. If the prescribed skin values are unset,
the coupling captures the initial surface-air values and holds them as the
surface reference.

There is no implemented multilayer soil temperature/moisture, vegetation or
prognostic land-water reservoir implied by that interface. Surface temperature
and moisture are prescribed/held forcing for the atmospheric exchange, not
the endpoint of an unshown soil energy equation. Radiation's surface absorption
receipt likewise does not automatically evolve skin temperature. Adding a
land model would require new state, an energy/water exchange contract, coupled
time integration, restart and admission rather than relabeling existing slab
controls.

\section{Controlled stochasticity, synthetic radar and diagnostics}

\subsection{Chaos is a perturbation owner}

The chaos registry provides disabled, initial-condition, boundary-layer and
full-stochastic families. These introduce controlled variability through
declared perturbation interfaces; \texttt{boundary\_layer} chaos does not
replace the selected physical PBL scheme. Keyed random streams, correlation
filters, member/seed identity and retained epochs belong to the experiment.
The \href{../src/chaos/base/random_generator.cpp}{random-stream implementation}
and \href{../src/chaos/runtime_state.cpp}{runtime state} prevent one logical
request or rejected update from silently becoming another request's history.

Moving-domain noise also has a declared Earth/frame geometry. Reusing an array
index after relocation is not a physical coordinate definition. Source
perturbations must preserve their underlying inventory/flux contract; scaling
a completed sedimentation divergence independently in every cell would not
preserve its surface precipitation exchange. Reproducibility claims remain
limited to the tested generator, arithmetic and continuation scope; a seed
alone is not a cross-library bitwise guarantee.

\subsection{Radar is a forward operator, not microphysics state}

The radar factory selects reflectivity, velocity and polarimetric ZDR operators;
product selection also exposes associated KDP and correlation outputs.
\href{../include/radar/radar_base.hpp}{RadarStateView} binds the chosen mass
and number inputs, dry density, pressure/theta, staggering, physical geometry
and Earth-frame radar location. Multiplying a dry-air mixing ratio by dry
density supplies species content per volume; treating it directly as volume
concentration would change the scattering calculation.

Radial velocity reconstructs face winds to the declared sample location before
line-of-sight projection and any configured scatterer fall-speed correction.
Terrain heights and grid origin matter to the beam direction. Reflectivity
and polarimetric tiers make particle/scattering assumptions beyond the mass
fields themselves; a microphysics sixth diameter moment is not automatically
every radar product. \href{../src/radar/runtime.cpp}{Radar runtime} constructs
only selected products and retains their operator state. These are diagnostics
of an accepted atmosphere, not additional forces feeding the dynamical state.

\subsection{Diagnostics supply contracts, analysis and receipts}

Diagnostics owns field membership/units, dry-mass and moisture inventories,
momentum and versioned energy, resolved derivatives, pressure decomposition,
water paths, storm tracking/rotation and parcel trajectories. Its diagnostic
pressure decomposition must not be confused with the rho--theta EOS pressure
commit. Its physical-domain sums must not include closure volumes merely
because those arrays contain finite values.

Some outputs require state across time: accumulated rainfall, parcel paths,
storm continuity and moving-window histories cannot be reconstructed from
one instantaneous field snapshot. Other products are derived lazily from an
accepted endpoint. \href{../src/diagnostics/derived/tranche_product_registry.cpp}{Product
registration} connects their dependencies to output selection, while
\href{../src/diagnostics/contracts/field_contract.cpp}{field contracts}
describe interpretation and bounds. Optional plotting products and required
scientific ledgers have different cadences and obligations.

An energy receipt attributes a change using a declared inventory; it does not
prove the corresponding physical heat or work law. Storm tracking is likewise
an algorithmic interpretation of fields with continuity and ambiguity rules,
not a uniquely identified physical object. The maintained
\href{../tests/diagnostics}{diagnostic tests} exercise those definitions and
their finite-coverage/failure conditions.

\section{Terrain, nesting and the common lifecycle}

The terrain factory provides flat, Bell and Sch\"ar geometry. Its physical
heights, interfaces, Jacobians and metric signs feed dynamics, transport,
column physics, diagnostics and output. Terrain changes numerical geometry;
it is not an actor transform applied only when drawing. Grid-attached and
Earth-fixed placement also have distinct motion semantics.

Nesting is coordinated within core infrastructure rather than by one independent
``nesting physics'' factory. Level runtimes own their schemes and contexts;
boundary fills, subcycling, restriction and reflux connect parent and child
inventories. Relocation/recycling additionally owns newly exposed regions,
private moments, held sources, physical origins and reports. The relevant
boundaries are \href{../src/core/infra/nested_execution.cpp}{nested execution},
\href{../src/core/infra/nested_scalar_feedback.cpp}{scalar feedback} and
\href{../src/core/infra/domain_recycling_state.cpp}{recycling state}.
Exact stationary and moving reference families are distinct from the broader
Engineering-incomplete moist/moving nesting item.

The \href{../src/core/runtime/headless_runtime.cpp}{headless runtime} brings
these modules together: initialization and admission create the state;
cadence owners refresh held physics; dynamics/orchestration advances it;
diagnostics and output describe committed endpoints; checkpointing retains the
full continuation state. Compute interfaces implement particular operations
inside that lifecycle, while utility helpers provide common parsing, cadence,
copy and registration rules. None replaces the higher-level physical owner.

Finally, \texttt{src/tools} contains table-generation and validation programs.
Generated coefficient/table files have explicit provisioning and identity
requirements. Their absence cannot be repaired by substituting a different
closure. These support tools, the included interfaces and the implementation
tree form a maintainable module map; long campaigns and independent review
remain separate evidence stages for the exact scientific combinations.
